\section{The rational boundary}\label{h3:app:rational-boundary}

We give the height-three argument of
\cite[Sections~2.5--2.6 and 3.9--6.3]{BSSW} for the rational algebra
used in \Cref{h3:thm:rational-compatibility}.  The separate trace
calculation in \Cref{h3:sec:rational} identifies the localization of
its degree-three generator.

\begin{proposition}[The rational height-three sphere]\label{h3:prop:rational-boundary}
For every prime \(p\geq5\), there is an isomorphism of graded algebras
\[
 \pi_*L_0T\simeq\Lambda_{\Q_p}(\alpha_1,\alpha_3,\alpha_5),
 \qquad |\alpha_i|=-i.
\]
\end{proposition}

The proof reduces this homotopy calculation to a comparison with
constant coefficients.  Put
\[
 \begin{aligned}
 \overline W&=W(\overline{\F}_p),&
 \breve K&=\overline W[1/p],\\
 \overline A&=\overline W[[u_1,u_2]],&
 \overline G&=P_3\rtimes\Gal(\overline{\F}_p/\F_p).
 \end{aligned}
\]
The map to be proved an isomorphism is the constants inclusion
\begin{equation}\label{h3:eq:rational-boundary-constants}
 H^*(\overline G,\overline W)[1/p]
 \longrightarrow H^*(\overline G,\overline A)[1/p].
\end{equation}
The compact Lie-algebra calculation identifies the source with the
exterior algebra in cohomological degrees $1,3,5$.  Comparisons of
the Lubin--Tate and Drinfeld towers identify the target with the
same graded vector space after adjoining to both sides an exterior
class $\epsilon$ of degree one.  A continuous equivariant additive
retraction onto $\overline W$ makes the source a direct summand.
Comparing dimensions, including the common $\epsilon$ factor,
then forces the complementary cohomology to vanish.  Since the
original constants inclusion is multiplicative, this proves the
algebra comparison.  Rational Morava descent converts it to the
asserted homotopy algebra.

The tower comparisons must be made integrally: the arithmetic
estimates give torsion bounds uniform over the charts and radii,
so that the derived limits and group invariants can be taken before
inverting $p$.  We first establish the Drinfeld model's actual
scalar equivalence.  Its building calculation also supplies the
properties used for the group $J_3$ in \Cref{h3:rat:trace-groups}.

\subsection{The Drinfeld model and its building}

\begin{lemma}[The reduced rank-three building]\label{h3:lem:rank-three-building}
Let $F$ be a non-Archimedean local field with normalized valuation $\nu$,
valuation ring $\mathcal O$, uniformizer $\pi$, and residue field $\kappa$
of cardinality $q$.  Let $\mathcal B$ be the reduced affine building of
$\mathrm{GL}_3(F)$, equivalently the building of $\mathrm{PGL}_3(F)$, and put
\[
 J_3(F)=\{g\in\mathrm{GL}_3(F):\nu(\det g)=0\}.
\]
Then $\mathcal B$ is a contractible, locally finite, two-dimensional
simplicial complex.  Each vertex has $2(q^2+q+1)$ adjacent vertices.
The group $J_3(F)$ preserves vertex types.  The setwise stabilizer of every
nonempty face is compact open in $J_3(F)$ and fixes that face pointwise.
Any closed chamber is a fundamental domain, including all its faces.
\end{lemma}

\begin{proof}
For $\mathrm{GL}_3$, the semisimple quotient used to define the reduced
building is the split group $\mathrm{PGL}_3$ of rank two.
The geometry statement in
\cite[Section~1.1.2, p.~6]{MeyerSolleveldBuildings}
applies over the local field $F$: the reduced building is locally finite,
has dimension the rank of this quotient, and is equivariantly contractible
for every compact subgroup of the quotient.  Taking the trivial compact
subgroup gives contractibility.

In the lattice-chain model of
\cite[Section~5.6, pp.~952--953]{DeligneFlicker2013}, vertices are homothety
classes $[L]$ of $\mathcal O$-lattices in $F^3$, and a chamber can be
represented by a strict cyclic chain
\[
 L_0\supset L_1\supset L_2\supset\pi L_0.
\]
Modulo $\pi L_0$ this is a complete flag in $\kappa^3$.  Thus a chamber
has three vertices and is a two-simplex.  An adjacent vertex to $[L]$ has
a unique representative $M$ with $\pi L\subsetneq M\subsetneq L$.
These representatives correspond to the nonzero proper subspaces of
$L/\pi L\cong\kappa^3$.  There are $q^2+q+1$ lines and the same number
of planes, giving the stated valency and a direct local-finiteness check.
Since the complex is locally finite, its metric polyhedral and weak
simplicial topologies agree locally on finite unions of simplices.  Its
realization is therefore a contractible CW complex.  The augmented
cellular chain complex over $\Z$ is exact, with free direct-sum terms
in degrees $0,1,2$.

Fix the standard lattice $\mathcal O^3$ and define
\[
 \operatorname{type}([g\mathcal O^3])=\nu(\det g)\pmod 3.
\]
Changing an integral basis changes the determinant by a unit; scalar
homothety changes its valuation by a multiple of three.  This type is
therefore well defined, and $J_3(F)$ preserves it.  The three vertices of
a chamber have all three types, since successive flag steps change the
lattice determinant valuation by one.

To send the standard chamber to a given chamber, start at the latter's
unique type-zero vertex.  Write its lattice as $M=g\mathcal O^3$ with
$\nu(\det g)=3m$, and replace $M$ by $M_0=\pi^{-m}M$.  Reanchor the
cyclic chain at $M_0$.  A basis $f_1,f_2,f_3$ of $M_0$ adapted to its
residue flag gives
\[
 M_i=\pi\mathcal O f_1+\cdots+\pi\mathcal O f_i
       +\mathcal O f_{i+1}+\cdots+\mathcal O f_3
 \quad(i=1,2).
\]
The map from the standard basis to $(f_1,f_2,f_3)$ has determinant
valuation zero, because its image lattice is $M_0$.  It lies in $J_3(F)$
and sends the standard chamber to the given chamber.

If $h\in J_3(F)$ stabilizes $[L]$, then $hL=aL$ for some $a\in F^\times$.
Taking determinant valuations relative to $L$ gives
$0=\nu(\det h)=3\nu(a)$, so $a$ is a unit and $hL=L$.  Hence the vertex
stabilizer is exactly $\operatorname{Aut}_{\mathcal O}(L)$, a conjugate
of the compact open group $\mathrm{GL}_3(\mathcal O)$.  A setwise face
stabilizer fixes each vertex because their types are distinct.  It is
therefore the finite intersection of these vertex stabilizers, hence is
compact open and fixes the face pointwise.

Every face lies in a chamber, and chamber transitivity sends it into the
standard chamber.  Its set of vertex types picks out a unique face there.
Type preservation also preserves the barycentric coordinates labeled by
those types, so no distinct points of the closed chamber are identified.
This proves the fundamental-domain assertion, including its faces.
\end{proof}

\begin{lemma}[The structure sheaf of the height-three Drinfeld model]
\label{h3:lem:drinfeld-structure-sheaf}
Let $\mathfrak H$ be the $p$-adic completion of the standard separated
strictly semistable weak formal model of $\Omega^2_{\Q_p}$ over $\Z_p$.
The canonical scalar map
\[
 \underline{\Z_p}\longrightarrow R\Gamma(\mathfrak H,\mathcal O)
\]
is an equivalence of condensed complexes, and is equivariant for the
action of $\GL_3(\Z_p)$.  Here the right side retains the derived
sections on $\mathfrak H\times\underline Q$ for every profinite set $Q$.
\end{lemma}

\begin{proof}
The imported model facts are given in
\cite[Sections~3.6,~3.8,~6.1--6.2,~6.4 and~7.1, and the proof of Proposition~6.5]{GKModel}.
After ordinary $p$-adic completion the local strictly semistable
equations and the special fiber are unchanged.  The model is separated:
each of the increasing weak-formal opens lies in a separated blowup
stage, and any two lie in one common such stage; the local closed-diagonal
criterion is preserved by completion.  The components and their
nonempty intersections are products of successive blowups of
projective spaces along rational linear subspaces.  They are indexed
by the vertices and simplices of the locally finite two-dimensional
Bruhat--Tits building $\mathcal B$ of \Cref{h3:lem:rank-three-building}.  Only the factors in dimensions
$0,1,2$ occur here.

We compute the structure-sheaf cohomology of exactly these factors.
For $\mathbf P^n_k$, $0\leq n\leq2$, over a field $k$, use the standard
affine cover.  On an intersection indexed by a nonempty set
$I\subset\{0,\ldots,n\}$, its degree-zero homogeneous Laurent
monomials have exponent vector $(a_0,\ldots,a_n)$ with
$\sum a_i=0$ and $a_i\geq0$ when $i\notin I$.  The Čech complex
decomposes by these monomials.  The constant monomial gives the
cochain complex of the full simplex and hence $k$ in degree zero.
For any other monomial its negative support
$N=\{i:a_i<0\}$ is nonempty and proper, and the monomial occurs
exactly for $I\supset N$.  Adding and deleting a fixed index outside
$N$, with the usual alternating sign, contracts this summand.
It follows that
\[
 R\Gamma(\mathbf P^n_k,\mathcal O)=k[0].
\]

The only nontrivial blowups needed for a surface factor are blowups
at the specified smooth rational points.  The calculation for one
such point is local on the surface and can be made in étale
coordinates $(u,v)$.  For the blowup of $\mathbf A^2_k$ at the
origin, the two standard charts and their overlap have rings
\[
 B_u=k[u,t],\qquad B_v=k[v,s],\qquad
 B_{uv}=k[u,t,t^{-1}],\qquad v=ut,\quad s=t^{-1}.
\]
In the common Laurent ring a monomial $u^a t^b$, with $a\geq0$ and
$b\in\Z$, belongs to $B_u$ when $b\geq0$ and to $B_v$ when $b\leq a$.
These two conditions cover all monomials.  Their intersection is
spanned by $0\leq b\leq a$, where
$u^a t^b=u^{a-b}v^b$.  Thus the affine Čech sequence
\[
 0\longrightarrow k[u,v]\longrightarrow B_u\oplus B_v
   \longrightarrow B_{uv}\longrightarrow0
\]
is exact.  The same sequence remains exact after localization and
flat étale base change, which are the local coordinate changes in
question.  The affine chart calculation therefore proves
$R\pi_*\mathcal O=\mathcal O$ for each required smooth-point blowup
$\pi$.  Blowing up an invertible ideal is an isomorphism.  This covers
the later strict transforms of rational lines on the surface and
the point centers on a projective line, all of which are Cartier
divisors.  Iterating the result, and using the finite affine-cover
double Čech complex for products over $k$, proves
\[
 R\Gamma(Y,\mathcal O_Y)=k[0]
\]
for every stratum $Y$ occurring in this height-three model.  This
argument makes no assertion about other rational varieties.
For a profinite parameter $Q$ the same calculations give
$C_{\hthreects}(Q,k)[0]$: continuous maps to each discrete section
module factor through a finite clopen partition of $Q$, and the
filtered union of the resulting finite powers preserves the exact
Čech sequences.

We next justify the resolution by intersections of components.
On a standard special-fiber chart its augmented form is the
alternating complex
\[
 k[t_0,\ldots,t_d]/(t_0\cdots t_r)
 \longrightarrow \bigoplus_i k[t_0,\ldots,t_d]/(t_i)
 \longrightarrow \bigoplus_{i<j} k[t_0,\ldots,t_d]/(t_i,t_j)
 \longrightarrow\cdots .
\]
For a nonzero monomial, let $Z$ be the nonempty set of indices
$i\leq r$ at which its exponent is zero.  Its coefficient complex
is the augmented cochain complex of the simplex on $Z$, so it is
exact.  This proves exactness monomial by monomial.  The calculation
is unchanged by the torus variables, localizations, and flat étale
base changes in the semistable charts.  Local finiteness of the
components now gives the global resolution by the products of their
intersection sheaves.  Taking derived global sections with a
profinite parameter and using the preceding stratum calculation
computes $R\Gamma(\mathfrak H_{\F_p}\times\underline Q,\mathcal O)$
by the three product terms below.  The arrow from
$C_{\hthreects}(Q,\F_p)$ is the separate canonical constants
augmentation:
\[
 C_{\hthreects}(Q,\F_p)\longrightarrow
 \prod_{\sigma\in\mathcal B_0}C_{\hthreects}(Q,\F_p)
 \longrightarrow
 \prod_{\sigma\in\mathcal B_1}C_{\hthreects}(Q,\F_p)
 \longrightarrow
 \prod_{\sigma\in\mathcal B_2}C_{\hthreects}(Q,\F_p),
\]
with the cellular incidence signs.  Products occur because the
strata are indexed by all cells of the building.

Here is a contraction valid for these product cochains.  By
\Cref{h3:lem:rank-three-building}, the augmented bounded cellular chain complex
$C_\bullet(\mathcal B;\Z)\to\Z$ is exact.  Its terms are free abelian,
and each cycle subgroup is free because it is a subgroup of a free
abelian group.  The short exact sequences of cycles therefore split.
Choosing splittings gives a chain contraction of the augmented
complex.  This contraction is not asserted to be equivariant.
Each basis cell maps to a finite cellular chain, since a cellular
chain group is a direct sum of its cells.  Dualizing into
$C_{\hthreects}(Q,\F_p)$ consequently gives a contraction of the
displayed product cochains: every output coordinate is a finite
linear combination of input coordinates and is continuous for the
product topology.  The contraction is natural in $Q$.

The canonical constants map into cellular cochains is equivariant
for the building action.  The contraction proves that its underlying
condensed map is an equivalence; the forgetful functor from equivariant
derived objects is conservative, so this same canonical map is an
equivariant equivalence.  In particular
\[
 \underline{\F_p}\xrightarrow{\ \simeq\ }
 R\Gamma(\mathfrak H_{\F_p},\mathcal O).
\]
Finally the formal structure sheaf is derived $p$-complete, and
derived global sections preserve its inverse limit of reductions.
The cone of the actual scalar map
$\underline{\Z_p}\to R\Gamma(\mathfrak H,\mathcal O)$ is therefore
derived $p$-complete.  Flatness of the semistable model identifies
its reduction modulo $p$ with the zero cone just computed.  Every
successive reduction modulo $p^\ell$ is then zero, and derived
$p$-completeness makes the cone itself zero.  This proves the
claimed equivariant equivalence.
\end{proof}

\subsection{The rational boundary calculation}

\begin{proof}[Proof of \Cref{h3:prop:rational-boundary}]
All cohomology complexes below retain their condensed parameters;
all inverse limits and group invariants are derived and taken
before inverting $p$.

\medskip\noindent\emph{1. Arithmetic bounds.} Let
$K$ be complete discretely valued of mixed characteristic $(0,p)$
with perfect residue field, and let $C$ be a completed algebraic
closure.  The scalar class obtained from the cyclotomic logarithm
gives a natural map
\[
 \mathcal O_K[\epsilon]\longrightarrow
 R\Gamma(G_K,\mathcal O_C),\qquad |\epsilon|=1,
 \quad \epsilon^2=0.
\]
Its cone is cohomologically nonnegative, and in each degree its
cohomology is killed by a power of $p$.  The required bounds can
be chosen uniformly after finite tame extensions of $K$.
For the finitely many nonzero twists $j=1,2$, the same assertion
of bounded torsion holds for $H^i(G_K,\mathcal O_C(-j))$, with
$H^0=0$.  Here is the part of the proof that gives uniformity,
\cite[Sections~4.2--4.4]{BSSW}.  Choose a finite cyclotomic stage $B/K$ far enough out that the
remaining $\Z_p$-tower is totally ramified and sufficiently ramified
in the sense of \cite[Definition~4.2.5 and Lemma~4.2.6]{BSSW}, generated
by $\sigma$, with $\chi(\sigma)$ a generator of $1+p^s\Z_p$,
$s\geq2$.  Write $M=\mathcal O_{\widehat K_\infty}$ and
$V=\ker(t:\widehat K_\infty\to B)$ for normalized trace.
The estimates in \cite[(4.2.9), (4.2.15), Lemmas~4.2.12 and
4.2.17, and the proof of Proposition~4.2.18]{BSSW} are
\[
 \|t\|\leq |p|^{-1/(p-1)},\qquad
 \|(\sigma-1)^{-1}\|_V
       \leq |p|^{-1-1/[p(p-1)]}.
\]
For existence of the inverse, at a finite cyclic level
$\sigma-1$ has kernel $B$ and image the trace-zero subspace.
Its inverse there is bounded by the second estimate, so the
inverses extend compatibly to the completion.  The estimate follows
by induction from normalized layer trace
$t_n=p^{-1}\operatorname{Tr}$ and
\[
 x-t_nx=p^{-1}\sum_{a=1}^{p-1}
       (1-\sigma^{a p^{n-1}})x,\qquad
 \|t_n\|\leq |p|^{-p^{-n}}.
\]
Set $M_0=\mathcal O_B\oplus(V\cap M)\subset M$.
The first norm bound gives $pM\subset M_0$, so $M/M_0$ is killed
by $p$.  For $j\in\{-2,-1,1,2\}$, put
$\lambda=\chi(\sigma)^{-j}$;
then $v_p(\lambda-1)=s$, since $p\geq5$.
The inequality
$|\lambda-1|\,\|(\sigma-1)^{-1}\|_V<1$ follows from
$s\geq2>1+1/[p(p-1)]$.  A Neumann series therefore inverts
$\sigma-\lambda$ on $V$ with the same bound
\cite[Lemma~4.4.1]{BSSW}, so its cokernel on $V\cap M$ is killed by $p^2$.
On $\mathcal O_B$ its cokernel is killed by $p^s$.
Thus $H^1(\langle\sigma\rangle,M_0(j))$ is killed by $p^s$,
while the invariants vanish.  The sequence $M_0\to M\to M/M_0$
gives a $p^{s+1}$ bound for $H^1(\langle\sigma\rangle,M(j))$.
The rational decomposition $B\oplus V$ also shows directly that
$M(j)$ has no invariants under the tail.
Restriction from $\Gal(K_\infty/K)$ to this open tail injects
$H^1$, because the tail has no invariants.  The cyclotomic fixed-field calculation and affinoid perfectoid
integral acyclicity compare this with $\mathcal O_C(j)$
\cite[Lemmas~4.3.2 and~4.3.5]{BSSW}.
Before cyclotomic invariants, the perfectoid-tail degree-zero
comparison has almost-zero cohomological error.  The maximal
ideal of the completed valuation ring is flat and idempotent,
so this error is annihilated by $p$ in the equivariant derived
category.  In the resulting triangle after invariants, the error
$Y$ has $H^0(Y)=0$ and $pH^i(Y)=0$ for $i>0$
\cite[Lemma~4.3.5 and Theorem~4.4.3]{BSSW}.
This gives the conservative bound $p^{s+2}$ in degree one;
the remaining degrees are bounded as in the same triangle.
The untwisted calculation retains the constant summand and its
cyclotomic-logarithm class, giving the displayed map
$\mathcal O_K[\epsilon]$ and its bounded-torsion cone.
Sufficient ramification persists under tame base change.  For a
finite tame extension $L/K$, let $e=e(LB/B)$ be the effective
tame ramification index after the chosen initial cyclotomic stage
$B$.  If $d_n$ is the different exponent of a $p$-layer above $B$,
then the base-changed exponent is
\[
 d'_n=e d_n-(e-1)(p-1).
\]
Consequently the same tail, $s$, and bounds work for all the
tame extensions in question.
Only existence of these tame-uniform bounds is used; the sharp
numerical constants in the preprint are unnecessary.

\medskip\noindent\emph{2. Semistable comparison.}
\par\noindent First apply this to a small affine semistable formal $\mathcal O_K$-scheme
$\mathfrak U=\operatorname{Spf}(R_0)$ of relative dimension $d\leq2$,
where small means admitting an \'etale morphism to a standard semistable chart.
Write $U$ for its generic fiber, $\mathfrak U_C$ for its completed base
change to $\mathcal O_C$, and $U_C$ for the geometric generic fiber. Use the divisorial log
structure and the valuation log structure on the base. The
semistable logarithmic comparison of
\cite[Sections~1.5--1.6 and Theorem~4.11]{CK}, extending the smooth
comparison of \cite[Theorem~8.3]{BMS}, applies on \'etale charts
\[
 \mathcal O_C\langle T_0,\ldots,T_r,
             T_{r+1}^{\pm1},\ldots,T_d^{\pm1}\rangle/
             (T_0\cdots T_r-a),\qquad 0\ne a\in\mathfrak m_C.
\]
They identify the cohomology of
$L\eta_{\zeta_p-1}R\nu_*\widehat{\mathcal O}^{+}$ with
$\Omega^j_{\mathfrak U_C,\log}\{-j\}$ for $0\leq j\leq d$.
No properness is required for these local statements.
The logarithmic differentials on $\mathfrak U$ are coherent and flat
over $\mathcal O_K$, and their higher cohomology on this affine
vanishes. Their completed base changes are the displayed
differentials on $\mathfrak U_C$.
The needed completed projection formula can be checked modulo
$p$: a flat module modulo $p$ is free over the artinian local
ring $\mathcal O_K/p$, and continuous cochains on a compact group
commute with these direct sums, since a map to a discrete direct
sum has finite image.  Derived $p$-completeness then proves the
projection formula before reduction.  For $j>0$, the normalized
map from the Breuil--Kisin twist to the Tate twist has cokernel
killed by $p$, so the preceding arithmetic bounds apply to each
of the finitely many positive graded pieces.

For the comparison with the original complex, put
$C_U=R\Gamma(U_{C,\mathrm{pro\acute et}},\widehat{\mathcal O}^{+})$,
$I=(\zeta_p-1)\subset\mathcal O_C$, and $A_U=L\eta_I C_U$.
The ideal $I$ is $G_K$-stable, although its displayed generator need not
be fixed.  On the semistable charts above, the perfectoid-cover edge map
of \cite[Section~3.3, Equation~(3.3.1)]{CK} has almost-zero cone.
Its source is the continuous cochain complex for
$\Delta\simeq\Z_p^d$, represented by a $d$-variable Koszul complex
\cite[Lemma~3.7]{CK}.  Thus $H^q(C_U)$ is almost zero for $q>d$.
Moreover, \cite[Theorems~3.9 and 4.11]{CK} identifies $A_U$ with a
complex concentrated in degrees $[0,d]$.  The degree-zero cohomology
of $C_U$ is $I$-torsion-free; indeed, it is the formal structure ring
\cite[Proposition~4.4]{CK}.

Set $B_U=\tau^{\leq d}C_U$.  Lemma~6.9 of \cite{BMS}, together
with compatibility of $L\eta_I$ with truncation, gives maps
\[
 a:A_U\longrightarrow B_U,
 \qquad b:I^{\otimes d}\otimes B_U\longrightarrow A_U
\]
whose composites are the natural ideal inclusions.  The invariant
element $p^d\in I^d$ gives an equivariant map
$\mathcal O_C\to I^{\otimes d}$; composing it with $b$ gives
$b_p:B_U\to A_U$ with $ab_p=p^d$ and $b_pa=p^d$.
Hence the kernel and cokernel of each $H^q(a)$ are killed by $p^d$,
and the long exact sequence gives
$p^{2d}H^q(\operatorname{cone}(a))=0$.
The triangle
\[
 \operatorname{cone}(a)\longrightarrow
 \operatorname{cone}(A_U\to C_U)\longrightarrow
 \tau^{>d}C_U
\]
then shows that $p^{2d+1}$ kills every cohomology group of the full
cone: the last term has almost-zero cohomology, which is killed by $p$.
The full cone is connective.  This argument concerns its cohomology
and does not assert that it carries a strict
$\Z/p^{2d+1}$-module structure.

After $G_K$-invariants, the first-quadrant spectral sequence has
terms $H^s(G_K,H^t(\operatorname{cone}(A_U\to C_U)))$.
In total degree $k$ its filtration has at most $k+1$ pieces, each
killed by $p^{2d+1}$.  Therefore $p^{(2d+1)(k+1)}$ kills the
cohomology in that degree.  These conservative bounds depend on
$d$ and $k$, and are independent of the chart and its tame extension.
The positive logarithmic graded pieces retain the arithmetic
bounds proved in Step~1. Tame descent preserves these bounds, by
unramified semilinear descent and averaging over tame inertia.
We obtain the natural comparison
\[
 \mathcal O(\mathfrak U)[\epsilon]\longrightarrow
 R\Gamma(U_{\mathrm{pro\acute et}},\widehat{\mathcal O}^{+})
\]
with nonnegative cone and degreewise bounds independent of the
chosen charts and their tame extensions. For a general semistable
affine, choose an \'etale hypercover refining the \v Cech nerve of
a small-chart cover, with each term a finite disjoint union of
small affine charts. Such finite covers of the matching objects
exist because the base is quasi-compact and quasi-separated.
Apply \'etale hyperdescent to the actual natural comparison map.
Its levelwise cones are connective and have the uniform degreewise
bounds just proved. In a fixed total degree the descent spectral
sequence has only finitely many cosimplicial and cohomological
positions, so the same bounded-torsion conclusion follows. Gluing next
by a separated affine formal cover likewise preserves a bound in
every total degree, since the relevant \v Cech diagonal is finite.  This proves the part of
\cite[Sections~5.3--5.6]{BSSW} needed here.

\medskip\noindent\emph{3. The two global calculations.} The Lubin--Tate space
$\mathrm{LT}_{\breve K}$ is the open two-dimensional ball.
Exhaust it by closed balls of radii $|p|^{1/\ell}$ for primes
$\ell\ne p$; each has a smooth model after a tame extension.
The two-index system of integral functions modulo $p^j$, indexed
by radius and $j$, is Mittag--Leffler.  This coefficient limit is
taken over the fixed base $\breve K$, whose uniformizer is $p$;
the tame extensions were used for the local comparisons.  More
generally, with uniformizer $\varpi$, for a fixed inner radius $r$
and reduction exponent $j$ take
$M=\lceil(j+1)/\log_{|\varpi|}(r)\rceil$.  Restriction from any
$s\geq r^{1/(j+1)}$, with source reduced modulo any
$\varpi^{j'}$ for $j'\geq j$, kills modulo $\varpi^j$ all terms of
total degree at least $M$, and the surviving coefficients are
integral.  Conversely every monomial of degree less than $M$ is
already integral on the larger ball.  The image is therefore the
fixed module spanned by those monomials.  Continuous functions from a
profinite parameter preserve this image statement by finite clopen
partitions.  A cofinal diagonal reduces the two indices to a
countable Mittag--Leffler tower.  Thus the derived two-index limit is
$\overline A$ in degree zero, as in
\cite[Lemmas~6.1.2--6.1.3]{BSSW}.  The preceding local comparisons
and the countable inverse-limit spectral sequence give
\[
 \overline A[\epsilon]\longrightarrow
 R\Gamma(\mathrm{LT}_{\breve K,\mathrm{pro\acute et}},
          \widehat{\mathcal O}^{+})
\]
with nonnegative, degreewise bounded-torsion cone: in any degree
the limit spectral sequence uses only a limit in that degree and
a $\lim^1$ from the preceding degree, both with bounds uniform
over the radii.  No uniform bound over all cohomological degrees is needed.

The other space is $\Omega^2_{\Q_p}$, with the formal model
$\mathfrak H$ of \Cref{h3:lem:drinfeld-structure-sheaf}.  That lemma
gives the actual equivariant scalar equivalence
$\underline{\Z_p}\simeq R\Gamma(\mathfrak H,\mathcal O)$, including
all profinite parameters.  It supplies the structure-sheaf assertion
needed from \cite[Theorem~6.2.1]{BSSW}.
The semistable comparison now gives
\[
 \Z_p[\epsilon]\longrightarrow
 R\Gamma(\Omega^2_{\Q_p,\mathrm{pro\acute et}},
          \widehat{\mathcal O}^{+})
\]
with the same degreewise bounded-torsion conclusion.

\medskip\noindent\emph{4. The tower comparison.} The bridge between these calculations is an equivalence of
\emph{families}, not merely a bijection of geometric points.
The relative Rapoport--Zink/shtuka comparison
\cite[Theorem~24.2.5 and its proof, pp.~227--229]{ScholzeWeinsteinBerkeley},
using the family comparison of its Theorem~22.3.1, identifies the
universal integral
Tate local system with that in a modification
\[
 0\longrightarrow \mathcal T\otimes_{\Z_p}\mathcal O
   \longrightarrow\mathcal E\longrightarrow i_*\mathcal L
   \longrightarrow0,
\]
where $\mathcal T=T_p(\mathcal X)$ is the actual integral Tate local
system of the universal $p$-divisible group $\mathcal X$, of rank
three, and $\mathcal L$ is a line at the
untilt divisor, and $\mathcal E$ is fiberwise $V_3$.
The relative basic-stratum frame torsors
\cite[Theorems~III.2.4 and III.4.5]{FarguesScholze} identify the two quotient
presentations of this moduli problem.  Honda endomorphisms are
already defined over $\F_{p^3}$; their canonical action and the
basic-stratum equivalence descend from the geometric residue
field as in \Cref{h3:geom:curve-inputs}.  After the height-component
normalization on the Lubin--Tate side, this gives
\[
 [\mathrm{LT}_{\breve K}^{\diamond}/\overline G]
 \simeq[\Omega^{2,\diamond}_{\Q_p}/\GL_3(\Z_p)].
\]
Pulling back the two atlases gives a common space with commuting
pro-\'etale torsor actions.  Their \v Cech nerves therefore compute
the same integral condensed cohomology object; this is the
actual descent argument behind \cite[Theorem~3.9.1 and (3.9.3)]{BSSW}.
Taking invariants of the comparison cones remains harmless:
they are nonnegative, so each diagonal of the group-cohomology
spectral sequence contains finitely many bounded-torsion rows.
After inverting $p$ we obtain
\[
 H^*(\overline G,\overline A)[1/p]\otimes\Q_p[\epsilon]
 \simeq H^*(\GL_3(\Z_p),\Q_p)\otimes\Q_p[\epsilon].
\]
The extended residue-field action is retained here.
Witt Artin--Schreier gives the continuous unramified calculation
over $\overline W$ in \cite[Lemma~3.8.5(1)]{BSSW};
its consequence for the full extended stabilizer is
\cite[Lemma~3.8.5(2)]{BSSW}.  Here $P_3$ fixes $\overline W$ and the
residue-field group acts by Witt Frobenius, exactly as in the chosen
Honda semilinear action.  The coefficient conventions also agree:
for every profinite $Q$, cochains on $Q\times\overline G^a$ with
values in the discrete $W_\ell(\overline{\F}_p)$ have finite image.
Lifting those finitely many values makes cochain reduction surjective.  Their derived inverse limit is therefore
the continuous $\overline W$ cochain complex.  Every continuous
cochain on this compact domain with values in $\overline W[1/p]$
has bounded image, so inverting $p$ gives the rational cochains
used in part~(2).

To pass to the prescribed finite field $k$, put
$\Delta_k=\Gal(\overline{\F}_p/k)$.  Its action fixes $u_1,u_2$
and acts by coefficient Frobenius.  Each quotient
$\overline A/(p,u_1,u_2)^a$ is a finite direct sum of truncated
Witt coefficient modules.  Witt Artin--Schreier therefore identifies
its $\Delta_k$-invariants with
$W(k)[[u_1,u_2]]/(p,u_1,u_2)^a$ and makes its positive
$\Delta_k$-cohomology vanish.  The coefficient reductions, including
cochains with profinite parameters, are surjective by finite-value
lifting.  Taking the derived inverse limit gives
\[
 R\Gamma(\Delta_k,\overline A)\simeq W(k)[[u_1,u_2]],
\]
equivariantly for $P_3$, since $k$ contains the Honda endomorphism
field.  Hochschild--Serre now identifies the algebraically closed
residue-field calculation with the finite-field one.  The trace-one
contraction of \Cref{h3:lem:framework-semilinear} supplies the finite
semilinear descent used by \Cref{h3:prop:exact-descent}.

\medskip\noindent\emph{5. Constants and rational homotopy.} The additive splitting in
\cite[Proposition~2.5.1]{BSSW} has a direct construction.
Compose each power operation with the $K(3)$-local transfer and
write the resulting operations as $\beta_m$.  Transfer transitivity
and the power-operation addition formula give
$\beta_m(a+b)=\sum_{i+j=m}\beta_i(a)\beta_j(b)$.
Thus $a\mapsto\sum_m\beta_m(a)t^m$, followed by reduction to
$\overline\F_p$ and projection from big to $p$-typical Witt vectors,
is an equivariant additive map
$\gamma:\overline A\to\overline W$.
For continuity, write the localized unit as the filtered colimit of its
compact Moore stages $M_k$. The spectrum
$L_{K(3)}\Sigma^\infty_+B\Sigma_m$ is dualizable for the finite
group $\Sigma_m$ \cite[Corollary~8.7]{HS99}; its tensor product with
the compact object $M_k$ is therefore compact. Compactness makes the map from
$M_k\otimes L_{K(3)}\Sigma^\infty_+ B\Sigma_m$ factor through some
finite tensor stage $(M_j^{\otimes m})_{h\Sigma_m}$; hence each
power operation is continuous for the Moore topology
\cite[Lemma~2.5.4]{BSSW}.
Transfers are maps of $K(3)$-local spectra and are continuous
for the same topology \cite[Proposition~11.1(b)]{HS99}.
For $\overline A=\pi_0\overline E$, where $\overline E$ is the
Morava theory used here, the Moore topology is the
$\mathfrak m=(p,u_1,u_2)$-adic topology by the following local check.
Choose the unit-compatible generalized Moore spectra $\mathbb S/J_j$
of \cite[Proposition~4.22]{HS99}.  Their $K(3)$-local duals
$M_j=D_{K(3)}L_{K(3)}(\mathbb S/J_j)$ are compact stages for the unit
by \cite[Corollary~7.11 and Theorem~8.5]{HS99}.  Duality carries restriction
along the dual unit exactly to the map induced by
$\mathbb S\to\mathbb S/J_j$.  The latter spectrum is finite, so
$\overline E\otimes(\mathbb S/J_j)$ remains $K(3)$-local.
The unit-normalized Landweber formula
\cite[Introduction and Example~0.1(d)]{HoveyStrickland} and the Moore
coefficient calculation \cite[Definition~4.12]{HS99} identify the restriction as
\[
 \overline A\longrightarrow[M_j,\overline E]
   \cong\pi_0(\overline E\otimes(\mathbb S/J_j))
   \cong\overline A/\overline J_j.
\]
Here $\overline J_j$ is the image of the Moore ideal after
removing invertible periodicity factors.  The normalized images of
$p,v_1,v_2$ form a regular sequence of length three in the
three-dimensional regular local ring $\overline A$.  Their ideal
is therefore $\mathfrak m$-primary.  The pure-power ideals
$\overline J_j$, whose exponents tend to infinity, are cofinal with
the powers of $\mathfrak m$,
which proves the claimed equality of topologies from the restriction
kernels.  Thus each composite $\beta_m:\overline A\to\overline A$
is $\mathfrak m$-adically continuous.  The coefficientwise
Witt-vector construction is therefore continuous.
Its restriction $f$ to $\overline W$ is the identity modulo $p$;
continuity and additivity make it $\Z_p$-linear.  Write
$f=\hthreeid+p h$ with a continuous $\Z_p$-linear endomorphism $h$.
The convergent series $\sum_{n\geq0}(-p h)^n$ is its inverse.
Hence $f^{-1}\gamma$ is an equivariant retraction and
$\overline A=\overline W\oplus\overline A^c$ additively.

The compact comparison used to compute the constants
has the following integral hypotheses.  Choose a deep normal uniform
open in $P_3$ and in $\GL_3(\Z_p)$.  At $p\geq5$ its lower-$p$-series
valuation is saturated and equi-$p$-valued with denominator $e=1$
and generators in degree one.  The trivial finite free module $\Z_p$
satisfies the hypotheses of \cite[Theorem~3.3.3]{HuberKingsNaumann}; its continuous
comparison is cup compatible.  The coefficient and restriction
naturality in \cite[Theorem~3.1.1(1)--(3)]{HuberKingsNaumann}
identifies its rationalization with Lazard's comparison
\cite[V.2.4.9--V.2.4.10]{Lazard} and makes
it equivariant for conjugation.  Rational Hochschild--Serre then
takes invariants under the finite quotient, whose higher cohomology
vanishes by averaging.  The full exterior calculation for the compact
matrix and stabilizer groups, together with the Lie-algebra
calculation and trivial adjoint action, is given by
\cite[Proposition~3.8.1 and Lemma~3.8.2]{BSSW}.
At height three its generator degrees are $1,3,5$.
The unramified and finite semilinear descent from Step~4, with
\cite[Lemma~3.8.5(2)]{BSSW} for the full extended stabilizer, therefore identifies both
\[
 H^*(\overline G,\overline W)[1/p]
 \quad\text{and}\quad H^*(\GL_3(\Z_p),\Q_p)
\]
with that exterior algebra over $\Q_p$.

The displayed tower comparison makes the cohomology of
$\overline A$ finite-dimensional in every degree.  Its additive
decomposition into constants and $\overline A^c$ is a decomposition
of continuous coefficient complexes.  If
$d_i=\dim_{\Q_p}H^i(\overline G,\overline A^c)[1/p]$, comparison
with the identical constant algebra, including the common factor
$\Q_p[\epsilon]$, gives $d_i+d_{i-1}=0$ (with $d_{-1}=0$).
Thus every $d_i$ is zero.
The constants inclusion is a ring map, so its resulting
cohomology isomorphism is multiplicative.  This proves
\eqref{h3:eq:rational-boundary-constants}.

Apply the local relative Morava descent of
\Cref{h3:prop:exact-descent} to $T$.  The central procyclic subgroup
generated by $1+p$ acts on internal degree $2t$ through
$(1+p)^{-t}$.  Its continuous two-term resolution has differential
$(1+p)^{-t}-1$ on that coefficient module.  For $t\ne0$ this scalar
is nonzero in $\Q_p$, so the rationalized two-term complex is
acyclic; Hochschild--Serre makes every nonzero internal row vanish.
The locally proved relative spectral sequence is strongly convergent
with finite cohomological amplitude.  Exact rationalization preserves
its bounded convergence.  Only the internal-degree-zero row remains, with
one filtration degree in each total degree.  Multiplicativity
therefore yields
\[
 \pi_*L_0T=\Lambda_{\Q_p}(\alpha_1,\alpha_3,\alpha_5),
 \qquad |\alpha_i|=-i.
\]
This establishes the rational boundary used in
\Cref{h3:thm:rational-compatibility}; the image of its chosen
primitive is determined by the separate map calculation in
\Cref{h3:sec:rational}.
\end{proof}
